% Generated draft: requires mathematical and source review.
\documentclass[11pt]{article}
\usepackage[margin=25mm]{geometry}
\usepackage{fontspec}
\setmainfont{Latin Modern Roman}
\usepackage{amsmath,amssymb,mathtools,unicode-math}
\setmathfont{Latin Modern Math}
\usepackage{enumitem,needspace,xurl,hyperref}
\setlength{\parindent}{0pt}
\setlength{\parskip}{5pt}
\setlength{\emergencystretch}{3em}
\newcommand{\N}{\mathbb{N}}
\newcommand{\Z}{\mathbb{Z}}
\newcommand{\Q}{\mathbb{Q}}
\newcommand{\R}{\mathbb{R}}
\newcommand{\C}{\mathbb{C}}
\newcommand{\F}{\mathbb{F}}
\newcommand{\A}{\mathbb{A}}
\newcommand{\PP}{\mathbb{P}}
\newcommand{\E}{\mathbb{E}}
\newcommand{\eps}{\varepsilon}
\newcommand{\dd}{\,\mathrm{d}}
\begin{document}
\section*{UnsolvedMath ID: 701393}
\textbf{Source problem code:} catalogue entry
\subsection{Definitions and mathematical statement}
For a group \(X\), define its derived series by \(X^{(0)}=X\) and \(X^{(i+1)}=[X^{(i)},X^{(i)}]\), where \([H,H]\) is the subgroup generated by all commutators of elements of \(H\). A group \(X\) is soluble if \(X^{(n)}=\{1\}\) for some integer \(n\geq 0\). Its derived length is
\[
d(X)=\min\{n\in\Z_{\geq 0}:X^{(n)}=\{1\}\}.
\]
In particular, \(d(\{1\})=0\). A maximal subgroup \(M\) of \(G\) means a proper subgroup with no subgroup strictly between \(M\) and \(G\).

\begin{enumerate}
\item[(a)] Does there exist an absolute constant \(k\in\Z_{\geq 0}\) such that, for every finite soluble group \(G\) and every maximal subgroup \(M<G\),
\[
d(G)-d(M)\leq k?
\]
\item[(b)] If such a constant exists, determine the least possible \(k\).
\end{enumerate}
Here \(k\) must be independent of both \(G\) and \(M\). Restricting \(k\) to nonnegative integers does not change the boundedness question, since the differences are nonnegative integers. The supplied record labels the problem open; this is a supplied status, not an independently verified literature assessment.
\subsection{Short English statement}
Is the difference between the derived length of a finite soluble group and that of any maximal subgroup bounded by a universal constant? If so, determine the smallest such constant.
\subsection{Sources}
\begin{itemize}
\item Local UnsolvedMath registry, permanent ID \(701393\), exact source code \texttt{catalogue entry}, title \emph{Kourovka Notebook Problem 17.91}. The supplied record labels the problem open and records an update date of 2026-10-16.
\item Attributed in the supplied record to V. S. Monakhov, proposed in 2010: \emph{The Kourovka Notebook (Unsolved Problems in Group Theory)}, 17th issue (2010), Problem 17.91. The background also identifies the 21st issue edition (2026), \texttt{arXiv:1401.0300v48}. The collection description names E. I. Khukhro and V. D. Mazurov as editors. These bibliographic details are supplied attribution and have not been independently verified.
\end{itemize}
\end{document}
